\section{Two quiet failures}
\label{sec:ch04-two-quiet-failures}
\index{projections!failure modes|(}

Both mistakes in this section were mine, and both were found by running the
service rather than by reading anything. Neither produces a stack trace, and one
of them produces no error at all, which is what earns them a section.

\subsection*{A field the schema promises and the resolver cannot keep}

\index{pagination!IncludeTotalCount}%
The first version of \code{browseProducts} did not carry the connection
attribute that turns a total count on:

\begin{minted}{csharp}
[UseConnection(IncludeTotalCount = true)]
\end{minted}

The schema it produced still had this in it, descriptions elided:

\begin{graphqlsdl}
type ProductConnection {
  edges: [ProductEdge!]
  nodes: [Product!]
  pageInfo: PageInfo!
  totalCount: Int! @cost(weight: "10")
}
\end{graphqlsdl}

A non-nullable \code{totalCount}, advertised to every client, on a field that
could not answer it:

\begin{minted}[fontsize=\footnotesize]{json}
{
  "errors": [
    {
      "message": "Cannot return null for non-nullable field.",
      "path": [ "browseProducts", "totalCount" ],
      "extensions": { "code": "HC0018" }
    }
  ],
  "data": null
}
\end{minted}

Note the \code{"data": null}. Non-null propagation walked all the way to the
root, so a client that asked for a page and a count got neither. Nothing about
this is subtle once it happens; the trouble is that the schema, which is the
thing clients read and generate code from, says the field is there and always
present.

Turning the option on fixes it, and costs less than I expected: the count
arrives as the correlated subquery in
section~\ref{sec:ch04-only-the-columns}, in the page's own statement. The
general lesson is the uncomfortable one. A connection type is generated for you,
and generated types can promise things your resolver has not been configured to
deliver. Export the SDL and read it, which chapter~\ref{ch:hotchocolate-quickly}
made a habit of for exactly this reason.

\subsection*{The identifier that was not there}

\index{projections!Parent requires}%
The second one produced no error whatsoever, and I nearly shipped it.

Ask \code{browseProducts} for identifiers and titles, and every product came
back with the same identifier:

\begin{minted}[fontsize=\footnotesize]{json}
{"nodes":[
  {"id":"00000000-0000-0000-0000-000000000000","title":"Beech Cutting Board"},
  {"id":"00000000-0000-0000-0000-000000000000","title":"Brant Ottoman"},
  {"id":"00000000-0000-0000-0000-000000000000","title":"Brant Two-Seat Sofa"}]}
\end{minted}

HTTP 200, no errors key, three products in the right order with the right
titles. Decode the \code{endCursor} that came with them and the zeros are in
there too, in the position where the sort tiebreaker should be. Page forward
from that cursor and you are seeking to a position that does not exist.

\index{projections!resolver-backed fields}%
The cause is a single line in Catalog:

\begin{minted}{csharp}
[ID]
public static Guid GetId([Parent] Product product) => product.Id;
\end{minted}

\code{Product.id} is produced by a resolver rather than read straight off the
property. I wrote it that way because \code{[ID]} is supposed to encode the raw
\code{Guid} into an opaque identifier, and
chapter~\ref{ch:schema-design-that-survives-change} discovers that at this tag
the attribute encodes nothing at all. What follows does not depend on that.
A field backed by a resolver breaks the projection whatever the resolver
does with the value. A projection builds its \code{SELECT} list out of the properties
named in the selection set. The selection set names \code{id}, which is a field
backed by a resolver, so the projection has nothing to select on its account,
so \code{products.id} is not in the statement, so the resolver is handed a
\code{Product} whose \code{Id} was never populated. \code{default(Guid)} is
sixteen zero bytes and no exception.

The fix is that \code{ParentAttribute} takes an optional string naming what the
resolver needs from its parent:

\begin{minted}{csharp}
[ID]
public static Guid GetId([Parent("Id")] Product product) => product.Id;
\end{minted}

Every cross-domain resolver on \code{Product} carries the same declaration now,
for the same reason: \code{price}, \code{availableQuantity}, \code{reviews} and
\code{averageRating} all key off an identifier that the selection set has no
reason to mention.

\subsection*{Why this one is the dangerous kind}

\index{projections!silent failure}%
Two details make this worse than an ordinary bug, and both are worth carrying
into any codebase that turns projection on.

The first is the shape of the failure. A missing column does not throw. It
produces a default value, and the default for the two most common key types is
zero or null: both legal, both serialisable, and both invisible to a test that
only counts elements.

The second is when it appears. Before the fix, \emph{not} selecting \code{id}
worked correctly, because the paging machinery adds the sort keys to the
selector on its own account, and \code{Id} is the tiebreaker. The failure only
showed up when a client asked for the field. So the version that goes wrong is
the version a real client sends, and the version in your smoke test, which asks
for a page and counts it, is the version that passes.

\subsection*{Verify it in Postman}

\index{Postman!projection assertions}%
Which is why the collection got three new requests for this chapter and one of
them exists solely to catch the zero.

The first asks for a page of three with identifiers, and asserts the page size,
the total count, the default ordering, and this:

\begin{minted}[fontsize=\footnotesize]{javascript}
pm.test("the projection kept the identifier", function () {
    page.nodes.forEach(function (node) {
        const raw = Buffer.from(node.id, "base64").toString("utf8");
        pm.expect(raw).to.not.include("00000000-0000-0000-0000-000000000000");
    });
});
\end{minted}

It decodes the opaque identifier and looks for the empty \code{Guid} inside it.
That is an ugly assertion and I would rather it did not have to exist, but a
correct-looking response is exactly what this failure produces, so the test has
to know what the wrong answer looks like.

One caveat, which belongs here rather than three chapters later. At this tag
\code{Product.id} answers with a raw \code{Guid}, because the attribute that
turns it into a base64 identifier is inert until
chapter~\ref{ch:schema-design-that-survives-change} switches on the machinery
behind it. The decode in that assertion is therefore handed text that was never
base64, and what comes back is binary noise that will not contain the printed
zero \code{Guid} whichever value the server sent. The assertion passes either
way. It could not have caught the bug it was written for, and it starts working
at the tag where identifiers really are base64 and the check reads the prefix
and the key bytes instead. The defect above is real and so is the fix. The
guard against it was not.

That request also stores its \code{endCursor} in a collection variable, and the
second reads it back and pages from it, asserting that the second page starts
where the first stopped. That is the cursor checked end to end rather than
eyeballed. The third sends a filter and an order together and asserts both took
effect.

\begin{minted}{text}
npx newman run postman/mosaic.postman_collection.json \
    -e postman/mosaic.local.postman_environment.json
\end{minted}

Ten requests, thirty-six assertions. \code{scripts/verify.ps1} runs the same
collection, and since this chapter it also starts PostgreSQL from
\code{docker-compose.yml}, waits for it to report healthy, and stops it again on
the way out. Both verification scripts now assert the lookup and statement
counts this chapter prints, so if a later chapter breaks the batching the script
says so before the prose does.
\index{projections!failure modes|)}
